\paragraph{单原子 Dirichlet 壳、residue 有限部与等功率 twisted 通道}
在单原子因子的周期--primitive 精确剥离、periodic residue 的 Fourier 反演以及 full/core residue 差的 rank-one 局域化已经建立之后，还可把同一原子在 Dirichlet 层继续压缩为一组完全闭式的后果。它们分别把 twisted 通道上的 ghost 贡献识别为精确 polylog 壳，把 residue 类上的 Dirichlet 级数写成 Hurwitz 壳并统一主极点残数，把 residue-wise 有限部常数的原子位移压缩为显式 digamma 项，并把 centered residue 扰动在全部非平凡字符上的单频能量锁定为同一常数。

\begin{theorem}[single atom 的 twisted Dirichlet 壳是精确的 polylog 壳]\label{thm:xi-time-part77-atomic-twisted-dirichlet-polylog-shell}
沿用定理 \ref{thm:xi-atomic-witt-cycle-primitive-exact-peeling} 的记号。对固定模数 \(m\ge 2\) 与扭曲通道 \(j\in \ZZ/m\ZZ\)，定义原子 periodic Dirichlet 级数
\[
\mathcal D^{\mathrm{cycle}}_{m,j}(s)
:=
\sum_{n\ge 1}a_n^{\mathrm{cycle}}(\omega_m^j)\,n^{-s}
\qquad (\Re s>1).
\]
则有完全闭式
\[
\mathcal D^{\mathrm{cycle}}_{m,j}(s)
=
2^{1-s}\operatorname{Li}_s(\omega_m^j),
\]
其中 \(j=0\) 时按 \(\operatorname{Li}_s(1)=\zeta(s)\) 解释。特别地：
\begin{enumerate}
  \item 当 \(j=0\) 时，
  \[
  \mathcal D^{\mathrm{cycle}}_{m,0}(s)=2^{1-s}\zeta(s),
  \]
  因而在 \(s=1\) 处具有简单极点；
  \item 当 \(1\le j\le m-1\) 时，\(\mathcal D^{\mathrm{cycle}}_{m,j}(s)\) 解析延拓为整个函数。
\end{enumerate}
\end{theorem}
\begin{proof}
由定理 \ref{thm:xi-atomic-witt-cycle-primitive-exact-peeling}，
\[
a_n^{\mathrm{cycle}}(u)=
\begin{cases}
2u^{n/2},& 2\mid n,\\
0,& 2\nmid n.
\end{cases}
\]
代入 \(u=\omega_m^j\) 得
\[
a_{2\ell}^{\mathrm{cycle}}(\omega_m^j)=2\omega_m^{j\ell},
\qquad
a_{2\ell+1}^{\mathrm{cycle}}(\omega_m^j)=0.
\]
因此
\[
\mathcal D^{\mathrm{cycle}}_{m,j}(s)
=
\sum_{\ell\ge 1}2\,\omega_m^{j\ell}(2\ell)^{-s}
=
2^{1-s}\sum_{\ell\ge 1}\frac{\omega_m^{j\ell}}{\ell^s}
=
2^{1-s}\operatorname{Li}_s(\omega_m^j),
\]
得到所示闭式。

当 \(j=0\) 时，上式退化为
\[
\mathcal D^{\mathrm{cycle}}_{m,0}(s)=2^{1-s}\zeta(s),
\]
故在 \(s=1\) 处有简单极点。

现设 \(j\neq 0\)。标准 root-of-unity 分解给出
\[
\operatorname{Li}_s(\omega_m^j)
=
m^{-s}\sum_{r=1}^{m}\omega_m^{jr}\,\zeta\!\left(s,\frac{r}{m}\right).
\]
对每个 \(r\)，Hurwitz zeta \(\zeta(s,r/m)\) 的唯一极点都位于 \(s=1\)，且留数同为 \(1\)。但
\[
\sum_{r=1}^{m}\omega_m^{jr}=0
\qquad (j\neq 0),
\]
故 \(s=1\) 处的极点严格相消，从而 \(\operatorname{Li}_s(\omega_m^j)\) 解析延拓为整个函数。乘上整因子 \(2^{1-s}\) 后，\(\mathcal D^{\mathrm{cycle}}_{m,j}(s)\) 亦为整个函数。证毕。
\end{proof}

\begin{theorem}[single atom 的 periodic residue Dirichlet 级数是 Hurwitz 壳，且全部 residue 主极点残数相等]\label{thm:xi-time-part77-atomic-residue-dirichlet-hurwitz-shell}
沿用定理 \ref{thm:xi-atomic-witt-energy-residue-l2-law} 与结论 \ref{con:xi-time-part57ba-atomic-residue-delta-flat-extremizer} 的记号。对固定模数 \(m\ge 2\) 与 residue 类 \(r\in \ZZ/m\ZZ\)，定义
\[
\mathcal A^{\mathrm{cycle}}_{m,r}(s)
:=
\sum_{n\ge 1}A^{\mathrm{cycle}}_{m,r}(n)\,n^{-s}
\qquad (\Re s>1).
\]
当 \(r\neq 0\) 时，下文把同一符号 \(r\) 也用于其唯一代表元 \(r\in\{1,\dots,m-1\}\)。
则成立：
\begin{enumerate}
  \item 对 \(r=0\)，有
  \[
  \mathcal A^{\mathrm{cycle}}_{m,0}(s)
  =
  2^{1-s}m^{-s}\zeta(s).
  \]
  \item 对 \(1\le r\le m-1\)，有
  \[
  \mathcal A^{\mathrm{cycle}}_{m,r}(s)
  =
  2^{1-s}m^{-s}\zeta\!\left(s,\frac{r}{m}\right).
  \]
\end{enumerate}
特别地，对每个 residue 类 \(r\in \ZZ/m\ZZ\)，\(\mathcal A^{\mathrm{cycle}}_{m,r}(s)\) 在 \(s=1\) 处都有简单极点，并且
\[
\operatorname*{Res}_{s=1}\mathcal A^{\mathrm{cycle}}_{m,r}(s)=\frac{1}{m}.
\]
此外，其 Laurent 常数项满足
\[
\mathcal A^{\mathrm{cycle}}_{m,0}(s)
=
\frac{1/m}{s-1}
+\frac{\gamma-\log(2m)}{m}
+O(s-1),
\]
\[
\mathcal A^{\mathrm{cycle}}_{m,r}(s)
=
\frac{1/m}{s-1}
+\frac{-\psi(r/m)-\log(2m)}{m}
+O(s-1)
\qquad (1\le r\le m-1),
\]
其中 \(\gamma\) 为 Euler 常数，\(\psi=\Gamma'/\Gamma\) 为 digamma 函数。
\end{theorem}
\begin{proof}
由结论 \ref{con:xi-time-part57ba-atomic-residue-delta-flat-extremizer}，
\[
A_{m,r}^{\mathrm{cycle}}(n)
=
2\,\mathbf 1_{\{2\mid n\}}\mathbf 1_{\{r\equiv n/2\;(\mathrm{mod}\,m)\}}.
\]
若 \(r=0\)，则非零项恰对应于 \(n=2mq\)（\(q\ge 1\)），故
\[
\mathcal A^{\mathrm{cycle}}_{m,0}(s)
=
2\sum_{q\ge 1}(2mq)^{-s}
=
2^{1-s}m^{-s}\zeta(s).
\]

若 \(1\le r\le m-1\)，则非零项恰对应于
\[
n=2(r+mq)\qquad (q\ge 0),
\]
故
\[
\mathcal A^{\mathrm{cycle}}_{m,r}(s)
=
2\sum_{q\ge 0}(2(r+mq))^{-s}
=
2^{1-s}m^{-s}\sum_{q\ge 0}\left(q+\frac{r}{m}\right)^{-s}
=
2^{1-s}m^{-s}\zeta\!\left(s,\frac{r}{m}\right).
\]
这就得到前两条公式。

已知 \(\zeta(s)\) 与全部 Hurwitz zeta \(\zeta(s,a)\)（\(0<a\le 1\)）在 \(s=1\) 的留数都等于 \(1\)。于是
\[
\operatorname*{Res}_{s=1}\mathcal A^{\mathrm{cycle}}_{m,r}(s)
=
\left.2^{1-s}m^{-s}\right|_{s=1}
=
\frac{1}{m}.
\]

再利用标准展开
\[
\zeta(s)=\frac{1}{s-1}+\gamma+O(s-1),
\qquad
\zeta(s,a)=\frac{1}{s-1}-\psi(a)+O(s-1),
\]
以及
\[
2^{1-s}m^{-s}
=
\frac{1}{m}\Bigl(1-(s-1)\log(2m)+O((s-1)^2)\Bigr),
\]
逐项相乘，即得所示 Laurent 常数项公式。证毕。
\end{proof}

\begin{corollary}[single atom 对 residue-wise 有限部常数只产生显式 digamma 位移]\label{cor:xi-time-part77-atomic-residue-finitepart-digamma-shift}
对固定模数 \(m\ge 2\) 与 residue 类 \(r\in \ZZ/m\ZZ\)，定义 full/core residue Dirichlet 级数
\[
\mathcal A^{\mathrm{full}}_{m,r}(s)
:=
\sum_{n\ge 1}A^{\mathrm{full}}_{m,r}(n)\,n^{-s},
\qquad
\mathcal A^{\mathrm{core}}_{m,r}(s)
:=
\sum_{n\ge 1}A^{\mathrm{core}}_{m,r}(n)\,n^{-s}.
\]
若把相应 Laurent 常数项记为
\[
\mathcal A^{\star}_{m,r}(s)
=
\frac{\kappa^{\star}_{m,r}}{s-1}
+\log\mathfrak M^{\star}_{m,r}
+O(s-1),
\qquad
\star\in\{\mathrm{full},\mathrm{core}\},
\]
并在 \(r\neq 0\) 时仍以 \(r\in\{1,\dots,m-1\}\) 表示其标准代表元。
则 single atom 对 residue-wise 有限部常数的修正完全显式：
\[
\log\mathfrak M^{\mathrm{full}}_{m,0}
-\log\mathfrak M^{\mathrm{core}}_{m,0}
=
\frac{\gamma-\log(2m)}{m},
\]
\[
\log\mathfrak M^{\mathrm{full}}_{m,r}
-\log\mathfrak M^{\mathrm{core}}_{m,r}
=
\frac{-\psi(r/m)-\log(2m)}{m}
\qquad (1\le r\le m-1).
\]
因此，single atom 在 residue-wise finite part 层并不与 core 混合，而只留下由 cyclotomic/Hurwitz 数据完全决定的一族 digamma 位移。
\end{corollary}
\begin{proof}
由定义，
\[
A^{\mathrm{full}}_{m,r}(n)=A^{\mathrm{core}}_{m,r}(n)+A^{\mathrm{cycle}}_{m,r}(n),
\]
故在 \(\Re s>1\) 上
\[
\mathcal A^{\mathrm{full}}_{m,r}(s)
=
\mathcal A^{\mathrm{core}}_{m,r}(s)+\mathcal A^{\mathrm{cycle}}_{m,r}(s).
\]
于是 \(\log\mathfrak M^{\mathrm{full}}_{m,r}-\log\mathfrak M^{\mathrm{core}}_{m,r}\) 恰等于 \(\mathcal A^{\mathrm{cycle}}_{m,r}(s)\) 在 \(s=1\) 处的 Laurent 常数项。再应用定理 \ref{thm:xi-time-part77-atomic-residue-dirichlet-hurwitz-shell} 的两条常数项公式，即得结论。证毕。
\end{proof}

\begin{corollary}[single atom 的 centered residue 扰动在全部非平凡 Fourier 模上等功率注入]\label{cor:xi-time-part77-atomic-centered-residue-equipower}
沿用定理 \ref{thm:xi-time-part65c-atomic-centered-residue-simplex} 的记号。对固定模数 \(m\ge 2\) 与 \(\ell\in \ZZ/m\ZZ\)，记 centered residue 扰动为
\[
\Delta^{(\ell)}_r:=2\,\mathbf 1_{\{r=\ell\}}-\frac{2}{m}
\qquad (r\in \ZZ/m\ZZ),
\]
并定义其归一化单频能量
\[
\mathcal E^{\mathrm{atom}}_{m,j}(\ell)
:=
\frac{1}{4}\left|\widehat{\Delta^{(\ell)}}(j)\right|^2
\qquad (1\le j\le m-1).
\]
则对所有 \(1\le j\le m-1\) 都有
\[
\mathcal E^{\mathrm{atom}}_{m,j}(\ell)=1,
\]
特别地，该值与 residue 位置 \(\ell\) 无关。换言之，single atom 在 \(\ZZ/m\ZZ\) 的全部非平凡字符通道上都以完全相同的功率注入，不存在任何优先频率。
\end{corollary}
\begin{proof}
定理 \ref{thm:xi-time-part65c-atomic-centered-residue-simplex} 已给出
\[
\widehat{\Delta^{(\ell)}}(0)=0,
\qquad
\widehat{\Delta^{(\ell)}}(j)=2\,\omega_m^{j\ell}
\qquad (1\le j\le m-1).
\]
因此对每个非平凡频率 \(j\) 都有
\[
\left|\widehat{\Delta^{(\ell)}}(j)\right|=2,
\]
从而
\[
\mathcal E^{\mathrm{atom}}_{m,j}(\ell)
=
\frac14\left|\widehat{\Delta^{(\ell)}}(j)\right|^2
=
\frac14\cdot 4
=
1.
\]
证毕。
\end{proof}

\noindent
这四条结果把单原子因子的 residue / twisted 接口进一步压缩到纯 Dirichlet 层的完全显式外壳：在 twisted 通道上，它恰等于一个根单位 polylog 壳；在 residue 通道上，它恰等于一族 Hurwitz 壳，并把全部主极点残数统一锁定为 \(1/m\)；相应地，full/core residue-wise finite part 常数之间的差也不再携带任何隐藏动力学信息，而只是一组显式的 digamma 位移；最后，centered residue 扰动在所有非平凡 Fourier 模上都具有同一功率，从而一切真正的频率选择性都只能来自 core，而不可能来自 single atom 本身。

\input{subsubsec__xi-time-protocol-conclusions-part78-phase-zeta-anomaly-splitting-classification}

\endinput
